\chapter{Introduction}
\label{chap:intro}

As organisations continue to grow, the need for a reliable and secure network infrastructure becomes increasingly important. CPUT requires an IT environment that can support its expanding operations while ensuring efficient management of users, devices, and network resources. To meet these requirements, a Windows Server 2019 environment was implemented using the hypervisor used during Semester~1.

This project focuses on the installation, configuration, and testing of key Windows Server services, including Active Directory Domain Services (AD DS), DNS, DHCP, Group Policy, Remote Access, and Windows Certificate Services. Screenshots were captured throughout the process to provide evidence of the installation and configuration steps performed.

\section{Background}
\label{sec:intro-background}
Windows Server 2019 is widely used in organisations to manage network resources, user accounts, security policies, and shared services from a central location. Services such as Active Directory, DNS, and DHCP help administrators manage users and devices more efficiently while ensuring that network communication functions correctly.

Security is also an important part of network administration. For this reason, several security policies were implemented in this project. These include disabling Command Prompt access, preventing the use of removable storage devices, restricting software installations, enforcing a minimum password length of ten characters, and requiring users to change their passwords regularly. These measures help reduce security risks and improve the overall protection of the network environment.

\section{Objectives}
\label{sec:intro-objectives}
The objectives of this report are:
\begin{itemize}
  \item To install and configure Windows Server 2019.
  \item To configure and test Active Directory Domain Services (AD DS).
  \item To configure and test DHCP services.
  \item To configure and test DNS services.
  \item To implement Group Policy settings that improve system security.
  \item To configure Remote Access services.
  \item To install and configure Windows Certificate Services.
  \item To verify that all configured services are functioning correctly.
  \item To document all installation, configuration, and testing procedures using screenshots.
\end{itemize}

\section{Report Structure}
\label{sec:intro-structure}
The remainder of this report is organised as follows:
\begin{itemize}
  \item Chapter~\ref{chap:install} covers the installation and initial configuration of Windows Server 2019. % Duy
  \item Chapter~\ref{chap:adds} discusses the configuration and testing of Active Directory Domain Services (AD DS). % Gian
  \item Chapter~\ref{chap:dhcp} explains the configuration and testing of DHCP services. % Peter
  \item Chapter~\ref{chap:dns} presents the configuration and testing of DNS services. % tbd
  \item Chapter~\ref{chap:gpo} describes the implementation of Group Policy settings. % Carl
  \item Chapter~\ref{chap:remote} discusses the configuration of Remote Access services. % Isaac
  \item Chapter~\ref{chap:certs} explains the installation and configuration of Windows Certificate Services. % Tihmo
  \item Chapter~\ref{chap:conclusion} provides the conclusion of the project. % Duy
  \item The final sections contain the bibliography and the appendix.
\end{itemize}
